\documentclass{article}
\usepackage{graphicx}
\usepackage{authblk}
\usepackage{indentfirst}
\usepackage{amssymb}
\usepackage{amsthm}
\usepackage{amsmath}
\usepackage{hyperref}
\usepackage{enumitem}
\usepackage{pgfplots}
\usepackage{tikz}
\usetikzlibrary{arrows.meta,calc}
\pgfplotsset{compat=1.18}
\newtheorem{theorem}{Theorem}
\newtheorem{corollary}{Corollary}[theorem]
\newtheorem{lemma}[theorem]{Lemma}
\newtheorem{definition}{Definition}
\newtheorem{problem}{Problem}
\newtheorem{solution}{Solution}
\newtheorem*{example}{Example}
\newtheorem{remark}{Remark}
\newtheorem{proposition}{Proposition}
\newtheorem{method}{Method}
\newtheorem{principle}{Principle}
\newcommand{\Span}{\operatorname{span}}
\newcommand{\R}{\mathbb{R}}
\newenvironment{amatrix}[1]{%
  \left(\begin{array}{@{}*{#1}{c}|c@{}}
}{%
  \end{array}\right)
}
\reversemarginpar
\hypersetup{
    colorlinks=true,
}
\begin{document}

%------- Title page -----------
\title{MATH 475: Honours Partial Differential Equations}
\author{W.R.YCH.}
\date{August \(31^{st}\), 2026}
\setcounter{Maxaffil}{0}
\renewcommand\Affilfont{\itshape\small}
\maketitle

%------- Abstract -----------
\noindent\rule{\textwidth}{0.4pt}
\thispagestyle{empty}
\begin{abstract}

\end{abstract}
\noindent\rule{\textwidth}{0.4pt}
\clearpage

%------- Table of Contents -----------
\thispagestyle{empty}
{
  \hypersetup{linkcolor=black}
  \tableofcontents
}
\clearpage

%------- Introduction -----------
\setcounter{page}{1}
\section{Introduction}
Taught by Adam Oberman. Mondays and Wednesdays from 14:35 to 15:55 in Stewart Biology Building N2/2. Course grading is divided as follows:
\begin{itemize}
    \item 10\% 5-6 assignments.
    \item 10\% attendance and participation.
    \item 10\% 3-5 quizzes.
    \item 20\% midterm.
    \item 50\% final exam.
\end{itemize}

% ============================================================
\section{Basic Definitions and Classification}
% ============================================================

\subsection{Ways to Write a PDE}
\marginpar{August 31, 2026}

Many of the classical equations are built from a single operator, the Laplacian
\begin{equation}
    \Delta u = \sum_{i=1}^{n} \frac{\partial^2 u}{\partial x_i^2}.
\end{equation}

\begin{center}
\begin{tabular}{lll}
    \textbf{Equation} & \textbf{Form} & \textbf{Type} \\
    \hline
    Heat & $u_t = \Delta u, \quad t \geq 0$ & parabolic \\
    Poisson & $\Delta u = f$ & elliptic (stationary) \\
    Laplace & $\Delta u = 0$ & elliptic (stationary) \\
    Wave & $u_{tt} = \Delta u$ & hyperbolic \\
\end{tabular}
\end{center}

Not every PDE fits this pattern. Nonlinear examples include Burgers' equation
\begin{equation}
    u_t + u u_x = 0,
\end{equation}
and the Navier-Stokes equations.

\subsection{Order and Solutions}
\marginpar{September 2, 2026}

\begin{definition}[Order]
    The order of a PDE is the order of the highest derivative appearing in it.
\end{definition}

For example $u_{xx} = 0$ is second order, and $u_t + u u_x = 0$ is first order.

\begin{definition}[Solution]
    Let the PDE be written as $F(u) = 0$ and let $k$ be its order. A function $u$ is a solution if
    \begin{enumerate}[label=(\arabic*)]
        \item $u \in C^k$, that is $u$ is $k$ times continuously differentiable, and
        \item $F(u) = 0$ at every point of the domain.
    \end{enumerate}
\end{definition}

\begin{example}[Transport equation]
    Let $u = u(x,t)$ and $F(u_t, u_x) = u_t + c u_x = 0$ with $c \in \mathbb{R}$. Let $\Phi$ be any differentiable function and set $u(x,t) = \Phi(x - ct)$. Then
    \begin{equation}
        u_x = \Phi'(x-ct), \qquad u_t = -c\,\Phi'(x-ct),
    \end{equation}
    so $u_t + c u_x = 0$. The whole family $\Phi$ solves the equation.
\end{example}

\begin{example}[Heat equation]
    Let $u_t = c^2 u_{xx}$ and take the ansatz $u(x,t) = e^{-c^2k^2 t}\sin(kx)$ for an arbitrary constant $k$. Then
    \begin{equation}
        u_t = -c^2k^2 u, \qquad
        u_x = k\,e^{-c^2k^2t}\cos(kx), \qquad
        u_{xx} = -k^2 e^{-c^2k^2t}\sin(kx) = -k^2 u,
    \end{equation}
    so $u_t = c^2 u_{xx}$.
\end{example}

\begin{example}[Wave equation]
    Let $u_{tt} = c^2 u_{xx}$ and take $u(x,t) = \cos(kct)\sin(kx)$. Then
    \begin{equation}
        u_{xx} = -k^2 u, \qquad u_{tt} = -c^2k^2 u,
    \end{equation}
    so $u_{tt} = c^2 u_{xx}$.
\end{example}

\subsection{Linear and Nonlinear}

\begin{definition}[Linear PDE]
    A PDE $F(u, u_x, x) = 0$ is linear if it can be rewritten as
    \begin{equation}
        L(u) = f(x),
    \end{equation}
    where $L$ is linear in $u$ and its derivatives and $f$ depends on $x$ only.
\end{definition}

\begin{example}
    Take $F(u, u_x, x) = 3u + u_x - x^2$. Setting $L(u, u_x) = 3u + u_x$ and $f(x) = x^2$ gives
    \begin{equation}
        F(u, u_x, x) = 0 \iff L(u, u_x) = f(x),
    \end{equation}
    so the equation is linear. The same definition works in higher dimensions.
\end{example}

\begin{definition}[Homogeneous]
    A linear PDE $L(u) = f(x)$ is homogeneous if $f \equiv 0$, and inhomogeneous otherwise.
\end{definition}

\begin{example}
    $u_t + a u_x = x^3$ is linear with $L(u) = u_t + a u_x$ and $f(x) = x^3$, so it is inhomogeneous.
\end{example}

\begin{example}
    $u_t + u u_x = f(x)$ is nonlinear, because the term $u u_x$ is not linear in $u$.
\end{example}

\subsection{Principle of Superposition}

\begin{principle}[Superposition]
    Let $L(u) = 0$ be linear and homogeneous. If $u_1$ and $u_2$ are solutions, then so is $u_1 + c\,u_2$ for every $c \in \mathbb{R}$.
\end{principle}

\begin{proof}
    By linearity of $L$,
    \begin{equation}
        L(u_1 + c u_2) = L(u_1) + c\,L(u_2) = 0 + c \cdot 0 = 0. \qedhere
    \end{equation}
\end{proof}

\begin{example}
    For $u_t = c^2 u_{xx}$, both
    \begin{equation}
        u_1(x,t) = e^{-c^2k^2t}\sin(kx)
        \qquad \text{and} \qquad
        u_2(x,t) = e^{-c^2 t}\sin(x)
    \end{equation}
    are solutions, so $u(x,t) = a_1 u_1(x,t) + a_2 u_2(x,t)$ is a solution for any $a_1, a_2 \in \mathbb{R}$.
\end{example}

\footnote{80\% chance that what we have written so far is not on the final.}

% ============================================================
\section{General Solutions}
% ============================================================

\subsection{Solving by Direct Integration}

\begin{example}
    Solve $u_x = 0$ for $u(x,y)$ on $\mathbb{R}^2$. The solution is $u(x,y) = f(y)$ for an arbitrary $f$.
\end{example}

\begin{example}
    Solve $u_{xx} = 0$. Set $v(x,y) = u_x$. Then $v_x = 0$, so by the previous example $v(x,y) = f(y)$, hence $u_x = f(y)$. Fixing $y$ and integrating in $x$ gives
    \begin{equation}
        u(x,y) = f(y)\,x + g(y).
    \end{equation}
    Check: $u_x = f(y)$ and $u_{xx} = 0$.
\end{example}

\begin{example}
    Solve $u_{xy} = 0$. Set $v(x,y) = u_x$. Then $v_y = 0$, so $v(x,y) = f(x)$, hence $u_x = f(x)$. Fixing $y$ and integrating in $x$ gives
    \begin{equation}
        u(x,y) = F(x) + g(y), \qquad F' = f.
    \end{equation}
    Check: $u_x = F'(x)$ depends on $x$ only, so $u_{xy} = 0$.
\end{example}

\begin{remark}
    In each case the general solution carries arbitrary \emph{functions} rather than arbitrary constants. An ODE of order $k$ has $k$ free constants; a PDE has free functions of the remaining variables.
\end{remark}

\subsection{Auxiliary Conditions}

Auxiliary conditions means boundary conditions. A typical problem on a domain $\Omega \subset \mathbb{R}^d$ reads
\begin{equation}
    \Delta u = f(x) \quad \text{in } \Omega,
    \qquad
    u(x) = g(x) \quad \text{on } \partial\Omega .
\end{equation}
These conditions are what cut the family of general solutions down to one.

\subsection{Operator Viewpoint}

Write a linear PDE as
\begin{equation}
    Lu = f,
\end{equation}
where $L$ is the operator, $u$ is the unknown and $f$ is the input. The linear algebra analogue is
\begin{equation}
    Mx = b, \qquad x, b \in \mathbb{R}^d, \qquad x = M^{-1}b,
\end{equation}
so solving a linear PDE is the infinite-dimensional version of inverting $M$.

% ============================================================
\section{Potential Theory}
% ============================================================

\subsection{Newton's Problem}

Place a point mass $M$ at the origin and a test mass $m$ at $x \in \mathbb{R}^3 \setminus \{0\}$. The gravitational force on $m$ has magnitude $GmM/\|x\|^2$ and points back toward the origin:
\begin{equation}
    \vec{F}(x) = -\frac{GmM}{\|x\|^2} \cdot \frac{x}{\|x\|} = -\frac{GmM}{\|x\|^3}\, x.
\end{equation}

\subsection{Toolkit}

\begin{lemma}[Derivative of distance is direction]
    Let $r(x) = \|x\|$ with $x = (x_1,x_2,x_3) \in \mathbb{R}^3$. Then for $x \neq 0$,
    \begin{equation}
        \frac{\partial r}{\partial x_i} = \frac{x_i}{r},
        \qquad
        \nabla r(x) = \frac{x}{r} = \hat{x}.
    \end{equation}
\end{lemma}

\begin{proof}
    Start from $r^2 = x_1^2 + x_2^2 + x_3^2$ and differentiate both sides with respect to $x_i$:
    \begin{equation}
        2r \frac{\partial r}{\partial x_i} = 2x_i
        \quad \Longrightarrow \quad
        \frac{\partial r}{\partial x_i} = \frac{x_i}{r}.
    \end{equation}
    This requires $r \neq 0$, so the identity fails at the origin.
\end{proof}

\begin{remark}
    $\nabla r$ is the unit vector pointing radially outward, so $\|\nabla r\| = 1$ everywhere it is defined.
\end{remark}

\subsection{Radial Chain Rule}

\begin{proposition}
    Let $f : (0,\infty) \to \mathbb{R}$ be differentiable and set $g(x) = f(r(x)) = f(\|x\|)$. Then for $x \neq 0$,
    \begin{equation}
        \nabla g(x) = f'(r)\, \nabla r(x) = f'(r)\, \frac{x}{r}.
    \end{equation}
\end{proposition}

\begin{proof}
    Componentwise, by the chain rule,
    \begin{equation}
        \frac{\partial g}{\partial x_i} = f'(r(x)) \frac{\partial r}{\partial x_i} = f'(r) \frac{x_i}{r}.
    \end{equation}
\end{proof}

So the gradient of any radial function points radially, and $f'$ alone controls the magnitude.

\begin{example}[Gradient refresher, non-radial]
    For $g(x) = x_1 x_2 + x_3^2$,
    \begin{equation}
        \nabla g(x) = \left( \frac{\partial g}{\partial x_1}, \frac{\partial g}{\partial x_2}, \frac{\partial g}{\partial x_3} \right) = (x_2,\; x_1,\; 2x_3).
    \end{equation}
\end{example}

\begin{remark}
    The chain rule result above, $\nabla g(x) = f'(r(x))\nabla r(x) = f'(r)\dfrac{x}{r}$, is the tool we use whenever a function depends only on distance from the origin.
\end{remark}

\subsection{The Laplacian}

Everything below works in $\mathbb{R}^d$ for any $d \geq 1$. Write $x = (x_1, \dots, x_d)$.

\subsubsection*{First form: sum of second partials}

\begin{definition}[Laplacian]
    \begin{equation}
        \Delta = \sum_{i=1}^{d} \frac{\partial^2}{\partial x_i^2},
        \qquad
        \Delta f(x) = \sum_{i=1}^{d} \frac{\partial^2 f}{\partial x_i^2}(x).
    \end{equation}
\end{definition}

In $d = 3$ this reads
\begin{equation}
    \Delta f(x) = \frac{\partial^2 f}{\partial x_1^2}(x) + \frac{\partial^2 f}{\partial x_2^2}(x) + \frac{\partial^2 f}{\partial x_3^2}(x).
\end{equation}

\begin{example}
    Let $f(x_1,x_2,x_3) = x_1^2 - x_2^2$. Then
    \begin{equation}
        \frac{\partial^2 f}{\partial x_1^2} = 2, \qquad
        \frac{\partial^2 f}{\partial x_2^2} = -2, \qquad
        \frac{\partial^2 f}{\partial x_3^2} = 0,
    \end{equation}
    so $\Delta f = 2 - 2 + 0 = 0$. A function with $\Delta f = 0$ is called \emph{harmonic}.
\end{example}

\subsubsection*{Second form: divergence of the gradient}

\begin{definition}[Divergence]
    For a vector field $\vec{v}(x) = (v^1(x), \dots, v^d(x))$, where the superscripts label components and are not powers,
    \begin{equation}
        \nabla \cdot \vec{v} = \operatorname{div} \vec{v} = \sum_{i=1}^{d} \frac{\partial v^i}{\partial x_i}(x).
    \end{equation}
\end{definition}

Applying this to the gradient field $\vec{v} = \nabla f$, whose $i$-th component is $\partial f / \partial x_i$, gives
\begin{equation}
    \nabla \cdot (\nabla f) = \sum_{i=1}^{d} \frac{\partial}{\partial x_i}\left( \frac{\partial f}{\partial x_i} \right) = \sum_{i=1}^{d} \frac{\partial^2 f}{\partial x_i^2} = \Delta f.
\end{equation}

\begin{proposition}
    \begin{equation}
        \Delta = \nabla \cdot \nabla = \operatorname{div} \circ \operatorname{grad}.
    \end{equation}
\end{proposition}

\begin{remark}
    This second form is the useful one for radial functions: combining it with the radial chain rule turns $\Delta$ into an ODE in $r$.
\end{remark}

\subsection{Harmonic Radial Functions}

\begin{proposition}[Radial Laplacian]
    For $g(x) = f(r)$, $r = \|x\|$ in $\mathbb{R}^d$, $x \neq 0$:
    \begin{equation}
        \Delta g = f''(r) + \frac{d-1}{r} f'(r).
    \end{equation}
\end{proposition}

Solve $\Delta f = 0$ in $d = 3$ with the ansatz $f(r) = r^a$:
\begin{equation}
    \Delta f = a(a-1)r^{a-2} + \frac{2}{r}\, a r^{a-1} = a(a+1)\, r^{a-2},
\end{equation}
which vanishes for $r > 0$ iff $a = 0$ or $a = -1$. So
\begin{equation}
    g(x) = A + \frac{B}{\|x\|}.
\end{equation}

\begin{remark}
    $1/\|x\|$ is the Newtonian potential: $\nabla(1/r) = -x/\|x\|^3$, the inverse-square field. Harmonic away from the origin only.
\end{remark}

% ============================================================
\section{First Order PDEs and the Solution Operator}
% ============================================================
\marginpar{September 9, 2026}

\subsection{The Underlying ODE}

Start from the ODE
\begin{equation}
    x'(s) = v(x(s), s),
\end{equation}
where $v$ is a velocity field. Assume $v(x,s)$ is smooth. By the existence and uniqueness theorem there is a unique trajectory passing through every point.

\subsection{The Solution Operator}

\begin{definition}[Solution operator]
    Given times $t_0$ and $t_1$, the solution operator is the map
    \begin{equation}
        S_{t_0 \to t_1}(x_0) = x(t_1),
    \end{equation}
    where $x(\cdot)$ solves the ODE with initial condition $x(t_0) = x_0$.
\end{definition}

\begin{proposition}[Identities]
    \begin{enumerate}
        \item $S_{t \to t}(x) = x$, that is $S_{t \to t} = \mathrm{Id}$.
        \item $S_{t_1 \to t_2} \circ S_{t_0 \to t_1} = S_{t_0 \to t_2}$.
        \item $S_{t_1 \to t_0} = \left( S_{t_0 \to t_1} \right)^{-1}$: running time backwards inverts the operator.
    \end{enumerate}
\end{proposition}

\subsection{Examples of Trajectories}

\begin{example}
    The curve $\vec{x}(s) = (x(s), y(s)) = (3e^{s-2}, 4e^{s-2})$ satisfies $\vec{x}'(s) = \vec{x}(s)$, so it is a trajectory of the velocity field $\vec{v}(x,y) = (x,y)$. It passes through $\vec{x}(2) = (3,4)$ and $\vec{x}(5) = (3e^3, 4e^3)$.
\end{example}

\begin{example}
    The curve $\vec{x}(s) = (3s+5, 5s+6)$ has constant velocity $\vec{v}(s) = \vec{x}'(s) = (3,5)$, with $\vec{x}(0) = (5,6)$ and $\vec{x}(2) = (11,16)$.
\end{example}

\begin{example}[Decoupled linear field]
    Take $\vec{v}(x,y) = (c_1 x, c_2 y)$ for constants $c_1, c_2$, so the ODE is
    \begin{equation}
        \begin{cases}
            x'(s) = c_1 x(s) \\
            y'(s) = c_2 y(s)
        \end{cases}
        \qquad
        \begin{cases}
            x(0) = x_0 \\
            y(0) = y_0.
        \end{cases}
    \end{equation}
    The solution is $x(s) = x_0 e^{c_1 s}$, $y(s) = y_0 e^{c_2 s}$. Check: $x'(s) = c_1 x_0 e^{c_1 s} = c_1 x(s)$, and similarly for $y$.
\end{example}

\begin{example}[Rotation field]
    Take $\vec{v}(x,y) = (-y, x)$. A circle of radius $r$, $\vec{x}(s) = (r\cos s, r\sin s)$, has $\vec{x}'(s) = (-r\sin s, r\cos s) = (-y, x)$, so circles about the origin are trajectories. The general solution of
    \begin{equation}
        \begin{cases}
            x' = -y \\
            y' = x
        \end{cases}
        \qquad
        \begin{cases}
            x(0) = x_0 \\
            y(0) = y_0
        \end{cases}
    \end{equation}
    is
    \begin{equation}
        x(s) = x_0 \cos s - y_0 \sin s, \qquad y(s) = x_0 \sin s + y_0 \cos s.
    \end{equation}
    The circle above is the case $x_0 = r$, $y_0 = 0$.
\end{example}

\subsection{Constant Coefficient Transport}

Consider
\begin{equation}
    a u_x + b u_y = 0, \qquad u = u(x,y), \qquad a, b \in \mathbb{R}.
\end{equation}

\begin{proposition}[Constancy along characteristics]
    Let $\vec{v} = (a,b)$ and let $\vec{x}(s) = (x(s), y(s))$ solve $\vec{x}'(s) = \vec{v}$, that is $x'(s) = a$ and $y'(s) = b$. Then $u(\vec{x}(s))$ is constant in $s$.
\end{proposition}

\begin{proof}
    By the chain rule,
    \begin{equation}
        \frac{d}{ds} u(x(s), y(s)) = u_x\, x'(s) + u_y\, y'(s) = a u_x + b u_y = 0. \qedhere
    \end{equation}
\end{proof}

The trajectories of $\vec{v}$ are called \emph{characteristics}. Since $u$ is constant along each one, its value is determined by wherever the characteristic meets the data.

\begin{proposition}[General solution]
    For any differentiable $f$,
    \begin{equation}
        u(x,y) = f(-bx + ay)
    \end{equation}
    solves $a u_x + b u_y = 0$.
\end{proposition}

\begin{proof}
    $u_x = -b\, f'(-bx+ay)$ and $u_y = a\, f'(-bx+ay)$, so $a u_x + b u_y = (-ab + ab) f'(-bx+ay) = 0$.
\end{proof}

\begin{remark}
    The argument $-bx + ay$ is constant exactly along the lines of direction $(a,b)$, which are the characteristics. When $b = 0$ the equation reduces to $a u_x = 0$ and the formula gives $u(x,y) = f(ay)$, a function of $y$ alone, as expected.
\end{remark}

\subsection{Adding Data}

\begin{example}
    Solve
    \begin{equation}
        \begin{cases}
            3u_x + 2u_y = 0 & \text{(PDE)} \\
            u(x,0) = x^3 & \text{(IC)},
        \end{cases}
    \end{equation}
    where the data is given on the line $\Gamma = \{(x,y) \mid y = 0\}$.
\end{example}

\begin{method}[Via the general solution]
    With $a = 3$, $b = 2$, the general solution is $u(x,y) = f(-2x + 3y)$. On $y = 0$ this reads $u(x,0) = f(-2x) = x^3$. Setting $z = -2x$, so $x = -z/2$, gives $f(z) = -z^3/8$. Therefore
    \begin{equation}
        u(x,y) = f(-2x + 3y) = -\frac{(-2x+3y)^3}{8} = \frac{(2x-3y)^3}{8}.
    \end{equation}
    Check: $u(x,0) = (2x)^3/8 = x^3$, and $3u_x + 2u_y = \tfrac{3}{8}\cdot 3(2x-3y)^2 \cdot 2 + \tfrac{2}{8}\cdot 3(2x-3y)^2 \cdot (-3) = 0$.
\end{method}

\begin{method}[Via characteristics]
    \begin{enumerate}[label=(\arabic*)]
        \item $u$ is constant along curves with $x'(s) = 3$, $y'(s) = 2$.
        \item These are the lines $x(s) = 3s + x_0$, $y(s) = 2s + y_0$.
        \item Given a target point $(x,y)$, find the point $(x_0, 0)$ on $\Gamma$ whose characteristic reaches it. From $y = 2s + 0$ we get $s = y/2$, and then $x = 3s + x_0$ gives $x_0 = x - 3y/2$.
        \item Since $u$ is constant along the characteristic, $u(x,y) = u(x_0, 0) = x_0^3$:
        \begin{equation}
            u(x,y) = \left( x - \frac{3y}{2} \right)^3 = \frac{(2x-3y)^3}{8}.
        \end{equation}
    \end{enumerate}
    Both methods agree.
\end{method}

% ============================================================
\section{Variable Coefficients and Inhomogeneous Transport}
% ============================================================
\marginpar{September 16, 2026}

\subsection{Problem Setup}

\begin{problem}[Linear transport with variable coefficients]
    Inputs:
    \begin{enumerate}[label=(\arabic*)]
        \item A curve $\Gamma \subset \mathbb{R}^2$ and data $u_0 : \Gamma \to \R$.
        \item Coefficients $a, b : \R^2 \to \R$.
    \end{enumerate}
    Find $u : \R^2 \to \R$ such that
    \begin{equation}
        \begin{cases}
            a(x,y)\,u_x + b(x,y)\,u_y = 0 & \text{(PDE)} \\
            u(x,y) = u_0(x,y) \quad \forall (x,y) \in \Gamma & \text{(IC)}.
        \end{cases}
    \end{equation}
\end{problem}

\subsection{Method of Characteristics}

\begin{method}[Homogeneous case]
    \begin{enumerate}[label=(\arabic*)]
        \item Solve the characteristic ODE
        \begin{equation}
            \begin{cases}
                x'(s) = a(x(s), y(s)) \\
                y'(s) = b(x(s), y(s))
            \end{cases}
            \qquad
            \begin{cases}
                x(0) = x_0 \\
                y(0) = y_0
            \end{cases}
            \qquad (x_0, y_0) \in \Gamma.
        \end{equation}
        \item Set $z(s) = u(x(s), y(s))$. By the chain rule
        \begin{equation}
            z'(s) = u_x\, x' + u_y\, y' = a u_x + b u_y = 0,
        \end{equation}
        so $z$ is constant and $z(s) = u_0(x_0, y_0)$.
        \item Given a target $(x,y)$, find $s_0$ and $(x_0, y_0) \in \Gamma$ with $x(s_0) = x$, $y(s_0) = y$. Then
        \begin{equation}
            u(x,y) = u_0\big(x_0(x,y),\, y_0(x,y)\big).
        \end{equation}
    \end{enumerate}
\end{method}

\subsection{Inhomogeneous Case}

Now allow a right hand side:
\begin{equation}
    a(x,y)\,u_x + b(x,y)\,u_y = c_1(x,y)\,u + c_2(x,y).
\end{equation}
The characteristics are unchanged, but $u$ is no longer constant along them. With $z(s) = u(x(s), y(s))$,
\begin{equation}
    z'(s) = a u_x + b u_y = c_1 z + c_2,
\end{equation}
so the method becomes a three-dimensional ODE system:
\begin{equation}
    \begin{cases}
        x' = a(x,y) \\
        y' = b(x,y) \\
        z' = c_1(x,y)\, z + c_2(x,y)
    \end{cases}
    \qquad
    \begin{cases}
        x(0) = x_0 \\
        y(0) = y_0 \\
        z(0) = u_0(x_0, y_0).
    \end{cases}
\end{equation}

\subsection{Worked Example}

\begin{example}
    Solve
    \begin{equation}
        \begin{cases}
            -y\,u_x + x\,u_y = u & \text{(PDE)} \\
            u(x,0) = g(x) & \text{(IC)}, \quad \Gamma = \{(x,0) \mid x > 0\}.
        \end{cases}
    \end{equation}
\end{example}

\begin{solution}
    Here $a = -y$, $b = x$, $c_1 = 1$, $c_2 = 0$.

    \textbf{Step 1: characteristic system.}
    \begin{equation}
        x' = -y, \qquad y' = x, \qquad z' = z,
        \qquad
        x(0) = x_0,\; y(0) = 0,\; z(0) = g(x_0).
    \end{equation}

    \textbf{Step 2: solve.} From the rotation field example, the general solution of $x' = -y$, $y' = x$ is
    \begin{equation}
        x(s) = a_1 \cos s - a_2 \sin s, \qquad y(s) = a_1 \sin s + a_2 \cos s.
    \end{equation}
    The initial conditions give $a_1 = x_0$, $a_2 = 0$, so
    \begin{equation}
        x(s) = x_0 \cos s, \qquad y(s) = x_0 \sin s, \qquad z(s) = g(x_0)\, e^{s}.
    \end{equation}

    \textbf{Step 3: invert.} Require $x(s_0) = x$, $y(s_0) = y$:
    \begin{equation}
        x = x_0 \cos s_0, \quad y = x_0 \sin s_0
        \quad \Longrightarrow \quad
        x_0 = \sqrt{x^2 + y^2}, \qquad s_0 = \arctan(y/x).
    \end{equation}

    \textbf{Step 4: read off $u$.}
    \begin{equation}
        u(x,y) = z(s_0) = g\!\left(\sqrt{x^2 + y^2}\right) \exp\!\big(\arctan(y/x)\big).
    \end{equation}
    In polar coordinates this is $u = g(r)\,e^{\theta}$, and the PDE is $\partial_\theta u = u$, which checks.
\end{solution}

\begin{remark}
    The characteristics are circles centred at the origin, so every point except the origin is reached from $\Gamma$. The formula with $\arctan$ is valid for $x > 0$; elsewhere use the polar angle $\theta \in [0, 2\pi)$ measured from the positive $x$-axis. Note $u$ jumps by a factor $e^{2\pi}$ across $\Gamma$ itself when $g \neq 0$, since going once around the circle multiplies $z$ by $e^{2\pi}$.
\end{remark}


\section{Misc}
$``\delta_p''$ point mass at $x=0, x\in\R$, $\int_{-\inf}^{\inf}\delta_0(x)dx=1$, $\delta_0(x)=0, x\neq 0$. Not possible for a function $\delta_0: \R\to\R$.
Not a function:
Definition: A function $f:\R\to\R$ is integrable if $\int_{-\inf}^{\inf}|f(x)|dx <\inf$. It is locally integrable if for all $a<b, a,b\in\R$ $\int_a^b|f(x)|dx <\inf$.
Example: $f(x)=1$, locally integrable $\int_a^b|f(x)|dx = (b-a)$, but not integrable $\int_{-\inf}^\inf 1dx = \inf$.
Example: suppose $f(x)=0,\forall |x|>R$, then integrable.
Example: $f(x) = \frac{1}{1+x^2}$ integrable but $f(x) = 1/x$, $\int_{0}^{1}1/xdx = logx |_0^1 = \inf$.

$f(x) = g(x), \forall x\neq3, f(3)\neq g(3), \int_{-\inf}^{\inf}f(x)dx = \int_{-\inf}^{\inf}g(x)dx$. Want to define $\delta_0$. Need for any continuous function $f(x)$, $\int_{-\inf}^{\inf}\delta_0(x)f(x)dx=f(0)$.

Distributiuons:
 -operational
 -theory

Heaviside:
$H(x) = \begin{cases}
    0, x< 0
    1, x\geq 0
\end{cases}$
Given $f(x)$ smooth
Define $H_{op}(f)=\int_{-\inf}^{+\inf}H(x)f(x)dx = \int_{0}^{\inf}f(x)dx$.

Given: $g:\R\to\R$
Define $g_{op}(f) = \int_{-\inf}^{\inf}g(x)f(x)dx$. $g_\epsilon = \begin{cases}
    0, x<0\\
    x/\epsilon, 0 \leq x \leq \epsilon\\
    1, x \geq \epsilon
\end{cases}$. Claim $\int_{-\inf}^{\inf}g_\epsilon'(x)f(x)dx$ is the average of $f$ over $[0\to\epsilon]$, $1/\epsilon\int_{0}^{\epsilon}f(x)dx$. btw $g_\epsilon'(x) = \begin{cases}
    0\\
    1/2\\
    0
\end{cases}$. f cts near 0: $\lim_{\epsilon\to 0} 1/\epsilon\int_{0}^{\epsilon}f(x)dx = f(0)$.

% ============================================================
\appendix
\section{Appendix}
% ============================================================

% ============================================================
\section{Solutions}
% ============================================================

\section{Useful Links}

\end{document}